
\documentclass{article}
\usepackage{colortbl}
\usepackage{makecell}
\usepackage{multirow}
\usepackage{supertabular}

\begin{document}

\newcounter{utterance}

\centering \large Interaction Transcript for game `clean\_up', experiment `0\_easy\_3obj\_en', episode 0 with gpt{-}5{-}mini{-}2025{-}08{-}07{-}t1.0.
\vspace{24pt}

{ \footnotesize  \setcounter{utterance}{1}
\setlength{\tabcolsep}{0pt}
\begin{supertabular}{c@{$\;$}|p{.15\linewidth}@{}p{.15\linewidth}p{.15\linewidth}p{.15\linewidth}p{.15\linewidth}p{.15\linewidth}}
    \# & $\;$A & \multicolumn{4}{c}{Game Master} & $\;\:$B\\
    \hline

    \theutterance \stepcounter{utterance}  
    & & \multicolumn{4}{p{0.6\linewidth}}{
        \cellcolor[rgb]{0.9,0.9,0.9}{
            \makecell[{{p{\linewidth}}}]{
                \texttt{\tiny{[P1$\langle$GM]}}
                \texttt{Sono il tuo game master, e stai giocando un gioco collaborativo con la seguente griglia come tavolo da gioco:} \\
\\ 
\texttt{```} \\
\texttt{ 1234567} \\
\texttt{╔═╤═══╤═╗} \\
\texttt{║◌│◌◌L│◌║ 1} \\
\texttt{╟─┼──┬┘◌║ 2} \\
\texttt{║◌│◌P│◌◌║ 3} \\
\texttt{╟─┴──┘◌◌║ 4} \\
\texttt{║◌◌◌◌◌C◌║ 5} \\
\texttt{║◌◌◌◌◌◌◌║ 6} \\
\texttt{║◌◌◌◌◌◌◌║ 7} \\
\texttt{╚═══════╝} \\
\\ 
\texttt{```} \\
\\ 
\texttt{* Il bordo superiore mostra le coordinate x che aumentano verso destra, e il bordo destro le coordinate y che aumentano verso il basso.} \\
\texttt{* I seguenti oggetti sono posizionati casualmente sulla tua griglia: 'C', 'L', 'P'.} \\
\\ 
\\ 
\texttt{L'altro giocatore vede una variante del tavolo da gioco, dove gli oggetti sono posizionati in diverse posizioni casuali. Non puoi vedere il tavolo dell'altro giocatore, e loro non possono vedere il tuo.} \\
\\ 
\\ 
\texttt{**Obiettivo:**} \\
\\ 
\texttt{Entrambi i giocatori devono spostare gli oggetti sui rispettivi sfondi in modo che oggetti identici finiscano alle stesse coordinate. Devi comunicare con l'altro giocatore per concordare una configurazione obiettivo comune.} \\
\\ 
\\ 
\texttt{**Regole:**} \\
\\ 
\texttt{* In ogni turno, puoi inviare esattamente uno dei seguenti due comandi:} \\
\texttt{ 1. `COMUNICA: <MESSAGE>`: per inviare un messaggio (tutto fino alla prossima interruzione di riga) all'altro giocatore. Lo inoltrerò al tuo partner.} \\
\texttt{2. `SPOSTA: <OBJECT>, (<X>, <Y>)`: per spostare un oggetto in una nuova posizione, dove <X> è la colonna e <Y> è la riga. Ti farò sapere se la tua mossa è andata a buon fine o no.} \\
\texttt{* Se non rispetti il formato, o invii diversi comandi contemporaneamente, devo penalizzarti.} \\
\texttt{* Se entrambi i giocatori accumulano più di 8 penalità, perdete entrambi il gioco.} \\
\texttt{* È vitale che comunichi con l'altro giocatore riguardo al tuo stato obiettivo! L'*unico* modo in cui puoi trasmettere la tua strategia all'altro giocatore è usando il comando `COMUNICA: <MESSAGE>`!} \\
\\ 
\\ 
\texttt{**Spostamento Oggetti**} \\
\\ 
\texttt{* Puoi spostare oggetti solo in celle entro i limiti della griglia. La cella di destinazione deve essere vuota, cioè deve contenere solo il simbolo '◌'.} \\
\texttt{* Se provi a spostare un oggetto in un posto che non è vuoto, o provi a spostarlo fuori dalla griglia, devo penalizzarti. Hai un altro tentativo.} \\
\texttt{* Prima di fare una mossa, controlla due volte che il posto di destinazione sia vuoto, e non contenga alcuna lettera, cornice, o linea!} \\
\\ 
\\ 
\texttt{**Fine del Gioco**} \\
\\ 
\texttt{Se pensi di aver raggiunto l'obiettivo di allineare tutti gli oggetti, puoi chiedere all'altro giocatore di finire il gioco inviando COMUNICA: finito?. Se l'altro giocatore ti chiede di finire il gioco, e rispondi con COMUNICA: finito!, the game will end.} \\
\\ 
\texttt{Entrambi i giocatori vincono se il gioco finisce entro 12 round, dove un round è definito come due giocatori che inviano ciascuno un comando valido.} \\
\\ 
\texttt{**Punteggio:**} \\
\\ 
\texttt{Più vicini sono gli oggetti identici in entrambi i tavoli da gioco, più punti ottieni. Le penalità riducono i tuoi punti. Riesci a battere il record?} \\
\\ 
\\ 
\texttt{Per favore invia un messaggio all'altro giocatore per iniziare il gioco!} \\
            }
        }
    }
    & \\ \\

    \theutterance \stepcounter{utterance}  
    & \multicolumn{4}{p{0.6\linewidth}}{
        \cellcolor[rgb]{0.9,0.9,0.9}{
            \makecell[{{p{\linewidth}}}]{
                \texttt{\tiny{[P1$\rangle$GM]}}
                \texttt{COMUNICA: Ciao! Propongo come configurazione obiettivo le mie posizioni attuali (tutte su celle libere): L(5,1), P(4,3), C(6,5). Puoi spostare i tuoi oggetti su queste coordinate? Se una di queste è occupata sul tuo lato, dimmi le tue posizioni attuali di C/L/P e proponi alternative. Io non muovo finché non confermi.} \\
            }
        }
    }
    & & \\ \\

    \theutterance \stepcounter{utterance}  
    & & & \multicolumn{2}{p{0.3\linewidth}}{
        \cellcolor[rgb]{0.9,0.9,0.9}{
            \makecell[{{p{\linewidth}}}]{
                \texttt{\tiny{[GM$|$GM]}}
                \texttt{COMUNICA: Ciao! Propongo come configurazione obiettivo le mie posizioni attuali (tutte su celle libere): L(5,1), P(4,3), C(6,5). Puoi spostare i tuoi oggetti su queste coordinate? Se una di queste è occupata sul tuo lato, dimmi le tue posizioni attuali di C/L/P e proponi alternative. Io non muovo finché non confermi.} \\
            }
        }
    }
    & & \\ \\

    \theutterance \stepcounter{utterance}  
    & & \multicolumn{4}{p{0.6\linewidth}}{
        \cellcolor[rgb]{0.9,0.9,0.9}{
            \makecell[{{p{\linewidth}}}]{
                \texttt{\tiny{[GM$\rangle$P2]}}
                \texttt{Sono il tuo game master, e stai giocando un gioco collaborativo con la seguente griglia come tavolo da gioco:} \\
\\ 
\texttt{```} \\
\texttt{ 1234567} \\
\texttt{╔═╤═══╤═╗} \\
\texttt{║◌│◌C◌│◌║ 1} \\
\texttt{╟─┼──┬┘◌║ 2} \\
\texttt{║◌│◌◌│◌◌║ 3} \\
\texttt{╟─┴──┘◌◌║ 4} \\
\texttt{║◌◌◌◌◌◌◌║ 5} \\
\texttt{║◌◌◌◌◌◌◌║ 6} \\
\texttt{║◌◌◌P◌◌L║ 7} \\
\texttt{╚═══════╝} \\
\\ 
\texttt{```} \\
\\ 
\texttt{* Il bordo superiore mostra le coordinate x che aumentano verso destra, e il bordo destro le coordinate y che aumentano verso il basso.} \\
\texttt{* I seguenti oggetti sono posizionati casualmente sulla tua griglia: 'C', 'L', 'P'.} \\
\\ 
\\ 
\texttt{L'altro giocatore vede una variante del tavolo da gioco, dove gli oggetti sono posizionati in diverse posizioni casuali. Non puoi vedere il tavolo dell'altro giocatore, e loro non possono vedere il tuo.} \\
\\ 
\\ 
\texttt{**Obiettivo:**} \\
\\ 
\texttt{Entrambi i giocatori devono spostare gli oggetti sui rispettivi sfondi in modo che oggetti identici finiscano alle stesse coordinate. Devi comunicare con l'altro giocatore per concordare una configurazione obiettivo comune.} \\
\\ 
\\ 
\texttt{**Regole:**} \\
\\ 
\texttt{* In ogni turno, puoi inviare esattamente uno dei seguenti due comandi:} \\
\texttt{ 1. `COMUNICA: <MESSAGE>`: per inviare un messaggio (tutto fino alla prossima interruzione di riga) all'altro giocatore. Lo inoltrerò al tuo partner.} \\
\texttt{2. `SPOSTA: <OBJECT>, (<X>, <Y>)`: per spostare un oggetto in una nuova posizione, dove <X> è la colonna e <Y> è la riga. Ti farò sapere se la tua mossa è andata a buon fine o no.} \\
\texttt{* Se non rispetti il formato, o invii diversi comandi contemporaneamente, devo penalizzarti.} \\
\texttt{* Se entrambi i giocatori accumulano più di 8 penalità, perdete entrambi il gioco.} \\
\texttt{* È vitale che comunichi con l'altro giocatore riguardo al tuo stato obiettivo! L'*unico* modo in cui puoi trasmettere la tua strategia all'altro giocatore è usando il comando `COMUNICA: <MESSAGE>`!} \\
\\ 
\\ 
\texttt{**Spostamento Oggetti**} \\
\\ 
\texttt{* Puoi spostare oggetti solo in celle entro i limiti della griglia. La cella di destinazione deve essere vuota, cioè deve contenere solo il simbolo '◌'.} \\
\texttt{* Se provi a spostare un oggetto in un posto che non è vuoto, o provi a spostarlo fuori dalla griglia, devo penalizzarti. Hai un altro tentativo.} \\
\texttt{* Prima di fare una mossa, controlla due volte che il posto di destinazione sia vuoto, e non contenga alcuna lettera, cornice, o linea!} \\
\\ 
\\ 
\texttt{**Fine del Gioco**} \\
\\ 
\texttt{Se pensi di aver raggiunto l'obiettivo di allineare tutti gli oggetti, puoi chiedere all'altro giocatore di finire il gioco inviando COMUNICA: finito?. Se l'altro giocatore ti chiede di finire il gioco, e rispondi con COMUNICA: finito!, the game will end.} \\
\\ 
\texttt{Entrambi i giocatori vincono se il gioco finisce entro 12 round, dove un round è definito come due giocatori che inviano ciascuno un comando valido.} \\
\\ 
\texttt{**Punteggio:**} \\
\\ 
\texttt{Più vicini sono gli oggetti identici in entrambi i tavoli da gioco, più punti ottieni. Le penalità riducono i tuoi punti. Riesci a battere il record?} \\
\\ 
\texttt{L'altro giocatore ha iniziato il gioco inviando questo messaggio:} \\
\\ 
\texttt{"Ciao! Propongo come configurazione obiettivo le mie posizioni attuali (tutte su celle libere): L(5,1), P(4,3), C(6,5). Puoi spostare i tuoi oggetti su queste coordinate? Se una di queste è occupata sul tuo lato, dimmi le tue posizioni attuali di C/L/P e proponi alternative. Io non muovo finché non confermi."} \\
\\ 
\texttt{Qual è il tuo primo comando?} \\
            }
        }
    }
    & \\ \\

    \theutterance \stepcounter{utterance}  
    & & & \multicolumn{4}{p{0.6\linewidth}}{
        \cellcolor[rgb]{0.9,0.9,0.9}{
            \makecell[{{p{\linewidth}}}]{
                \texttt{\tiny{[GM$\langle$P2]}}
                \texttt{COMUNICA: Confermo il tuo obiettivo: L(5,1), P(4,3), C(6,5) sono libere sul mio tavolo. Le mie posizioni attuali: C(4,1), P(4,7), L(7,7). Procedo a muovere in quest’ordine: P→(4,3), poi L→(5,1), poi C→(6,5). Non muovere nulla finché non ti dico di aver finito.} \\
            }
        }
    }
     \\ \\

    \theutterance \stepcounter{utterance}  
    & & & \multicolumn{2}{p{0.3\linewidth}}{
        \cellcolor[rgb]{0.9,0.9,0.9}{
            \makecell[{{p{\linewidth}}}]{
                \texttt{\tiny{[GM$|$GM]}}
                \texttt{COMUNICA: Confermo il tuo obiettivo: L(5,1), P(4,3), C(6,5) sono libere sul mio tavolo. Le mie posizioni attuali: C(4,1), P(4,7), L(7,7). Procedo a muovere in quest’ordine: P→(4,3), poi L→(5,1), poi C→(6,5). Non muovere nulla finché non ti dico di aver finito.} \\
            }
        }
    }
    & & \\ \\

    \theutterance \stepcounter{utterance}  
    & & \multicolumn{4}{p{0.6\linewidth}}{
        \cellcolor[rgb]{0.9,0.9,0.9}{
            \makecell[{{p{\linewidth}}}]{
                \texttt{\tiny{[P1$\langle$GM]}}
                \texttt{Il tuo messaggio è stato inoltrato all'altro giocatore.} \\
\\ 
\\ 
\texttt{Stai attualmente giocando il round 1 di massimo 12.} \\
\texttt{Avete accumulato collettivamente 0 di 8 penalità.} \\
\texttt{L'altro giocatore ha inviato questo messaggio:} \\
\\ 
\\ 
\texttt{"Confermo il tuo obiettivo: L(5,1), P(4,3), C(6,5) sono libere sul mio tavolo. Le mie posizioni attuali: C(4,1), P(4,7), L(7,7). Procedo a muovere in quest’ordine: P→(4,3), poi L→(5,1), poi C→(6,5). Non muovere nulla finché non ti dico di aver finito."} \\
\\ 
\\ 
\texttt{Qual è il tuo prossimo comando?} \\
            }
        }
    }
    & \\ \\

    \theutterance \stepcounter{utterance}  
    & \multicolumn{4}{p{0.6\linewidth}}{
        \cellcolor[rgb]{0.9,0.9,0.9}{
            \makecell[{{p{\linewidth}}}]{
                \texttt{\tiny{[P1$\rangle$GM]}}
                \texttt{COMUNICA: Perfetto, resto fermo. Quando hai completato: P→(4,3), L→(5,1), C→(6,5), inviami "finito?" e risponderò "finito!". Se incontri problemi (cella occupata o blocchi), avvisami e ricalibriamo.} \\
            }
        }
    }
    & & \\ \\

    \theutterance \stepcounter{utterance}  
    & & & \multicolumn{2}{p{0.3\linewidth}}{
        \cellcolor[rgb]{0.9,0.9,0.9}{
            \makecell[{{p{\linewidth}}}]{
                \texttt{\tiny{[GM$|$GM]}}
                \texttt{COMUNICA: Perfetto, resto fermo. Quando hai completato: P→(4,3), L→(5,1), C→(6,5), inviami "finito?" e risponderò "finito!". Se incontri problemi (cella occupata o blocchi), avvisami e ricalibriamo.} \\
            }
        }
    }
    & & \\ \\

    \theutterance \stepcounter{utterance}  
    & & \multicolumn{4}{p{0.6\linewidth}}{
        \cellcolor[rgb]{0.9,0.9,0.9}{
            \makecell[{{p{\linewidth}}}]{
                \texttt{\tiny{[GM$\rangle$P2]}}
                \texttt{Il tuo messaggio è stato inoltrato all'altro giocatore.} \\
\\ 
\\ 
\texttt{Stai attualmente giocando il round 2 di massimo 12.} \\
\texttt{Avete accumulato collettivamente 0 di 8 penalità.} \\
\texttt{L'altro giocatore ha inviato questo messaggio:} \\
\\ 
\\ 
\texttt{"Perfetto, resto fermo. Quando hai completato: P→(4,3), L→(5,1), C→(6,5), inviami "finito?" e risponderò "finito!". Se incontri problemi (cella occupata o blocchi), avvisami e ricalibriamo."} \\
\\ 
\\ 
\texttt{Qual è il tuo prossimo comando?} \\
            }
        }
    }
    & \\ \\

    \theutterance \stepcounter{utterance}  
    & & & \multicolumn{4}{p{0.6\linewidth}}{
        \cellcolor[rgb]{0.9,0.9,0.9}{
            \makecell[{{p{\linewidth}}}]{
                \texttt{\tiny{[GM$\langle$P2]}}
                \texttt{SPOSTA: P, (4, 3)} \\
            }
        }
    }
     \\ \\

    \theutterance \stepcounter{utterance}  
    & & & \multicolumn{2}{p{0.3\linewidth}}{
        \cellcolor[rgb]{0.9,0.9,0.9}{
            \makecell[{{p{\linewidth}}}]{
                \texttt{\tiny{[GM$|$GM]}}
                \texttt{SPOSTA: P, (4, 3)} \\
            }
        }
    }
    & & \\ \\

    \theutterance \stepcounter{utterance}  
    & & & \multicolumn{2}{p{0.3\linewidth}}{
        \cellcolor[rgb]{0.9,0.9,0.9}{
            \makecell[{{p{\linewidth}}}]{
                \texttt{\tiny{[GM$|$GM]}}
                \texttt{Spostato 'P' in (4,3) con successo. La tua griglia aggiornata appare così:} \\
\texttt{```} \\
\texttt{ 1234567} \\
\texttt{╔═╤═══╤═╗} \\
\texttt{║◌│◌C◌│◌║ 1} \\
\texttt{╟─┼──┬┘◌║ 2} \\
\texttt{║◌│◌P│◌◌║ 3} \\
\texttt{╟─┴──┘◌◌║ 4} \\
\texttt{║◌◌◌◌◌◌◌║ 5} \\
\texttt{║◌◌◌◌◌◌◌║ 6} \\
\texttt{║◌◌◌◌◌◌L║ 7} \\
\texttt{╚═══════╝} \\
\texttt{```} \\
            }
        }
    }
    & & \\ \\

    \theutterance \stepcounter{utterance}  
    & & \multicolumn{4}{p{0.6\linewidth}}{
        \cellcolor[rgb]{0.9,0.9,0.9}{
            \makecell[{{p{\linewidth}}}]{
                \texttt{\tiny{[P1$\langle$GM]}}
                \texttt{Il tuo messaggio è stato inoltrato all'altro giocatore.} \\
\\ 
\\ 
\texttt{Stai attualmente giocando il round 2 di massimo 12.} \\
\texttt{Avete accumulato collettivamente 0 di 8 penalità.} \\
\texttt{L'altro giocatore ha spostato un oggetto sulla sua griglia.} \\
\\ 
\\ 
\texttt{Qual è il tuo prossimo comando?} \\
            }
        }
    }
    & \\ \\

    \theutterance \stepcounter{utterance}  
    & \multicolumn{4}{p{0.6\linewidth}}{
        \cellcolor[rgb]{0.9,0.9,0.9}{
            \makecell[{{p{\linewidth}}}]{
                \texttt{\tiny{[P1$\rangle$GM]}}
                \texttt{COMUNICA: Ok, prosegui con le mosse; resto fermo e attendo il tuo "finito?".} \\
            }
        }
    }
    & & \\ \\

    \theutterance \stepcounter{utterance}  
    & & & \multicolumn{2}{p{0.3\linewidth}}{
        \cellcolor[rgb]{0.9,0.9,0.9}{
            \makecell[{{p{\linewidth}}}]{
                \texttt{\tiny{[GM$|$GM]}}
                \texttt{COMUNICA: Ok, prosegui con le mosse; resto fermo e attendo il tuo "finito?".} \\
            }
        }
    }
    & & \\ \\

    \theutterance \stepcounter{utterance}  
    & & \multicolumn{4}{p{0.6\linewidth}}{
        \cellcolor[rgb]{0.9,0.9,0.9}{
            \makecell[{{p{\linewidth}}}]{
                \texttt{\tiny{[GM$\rangle$P2]}}
                \texttt{Spostato 'P' in (4,3) con successo. La tua griglia aggiornata appare così:} \\
\texttt{```} \\
\texttt{ 1234567} \\
\texttt{╔═╤═══╤═╗} \\
\texttt{║◌│◌C◌│◌║ 1} \\
\texttt{╟─┼──┬┘◌║ 2} \\
\texttt{║◌│◌P│◌◌║ 3} \\
\texttt{╟─┴──┘◌◌║ 4} \\
\texttt{║◌◌◌◌◌◌◌║ 5} \\
\texttt{║◌◌◌◌◌◌◌║ 6} \\
\texttt{║◌◌◌◌◌◌L║ 7} \\
\texttt{╚═══════╝} \\
\texttt{```} \\
\texttt{Stai attualmente giocando il round 3 di massimo 12.} \\
\texttt{Avete accumulato collettivamente 0 di 8 penalità.} \\
\texttt{L'altro giocatore ha inviato questo messaggio:} \\
\\ 
\\ 
\texttt{"Ok, prosegui con le mosse; resto fermo e attendo il tuo "finito?"."} \\
\\ 
\\ 
\texttt{Qual è il tuo prossimo comando?} \\
            }
        }
    }
    & \\ \\

    \theutterance \stepcounter{utterance}  
    & & & \multicolumn{4}{p{0.6\linewidth}}{
        \cellcolor[rgb]{0.9,0.9,0.9}{
            \makecell[{{p{\linewidth}}}]{
                \texttt{\tiny{[GM$\langle$P2]}}
                \texttt{SPOSTA: L, (5, 1)} \\
            }
        }
    }
     \\ \\

    \theutterance \stepcounter{utterance}  
    & & & \multicolumn{2}{p{0.3\linewidth}}{
        \cellcolor[rgb]{0.9,0.9,0.9}{
            \makecell[{{p{\linewidth}}}]{
                \texttt{\tiny{[GM$|$GM]}}
                \texttt{SPOSTA: L, (5, 1)} \\
            }
        }
    }
    & & \\ \\

    \theutterance \stepcounter{utterance}  
    & & & \multicolumn{2}{p{0.3\linewidth}}{
        \cellcolor[rgb]{0.9,0.9,0.9}{
            \makecell[{{p{\linewidth}}}]{
                \texttt{\tiny{[GM$|$GM]}}
                \texttt{Spostato 'L' in (5,1) con successo. La tua griglia aggiornata appare così:} \\
\texttt{```} \\
\texttt{ 1234567} \\
\texttt{╔═╤═══╤═╗} \\
\texttt{║◌│◌CL│◌║ 1} \\
\texttt{╟─┼──┬┘◌║ 2} \\
\texttt{║◌│◌P│◌◌║ 3} \\
\texttt{╟─┴──┘◌◌║ 4} \\
\texttt{║◌◌◌◌◌◌◌║ 5} \\
\texttt{║◌◌◌◌◌◌◌║ 6} \\
\texttt{║◌◌◌◌◌◌◌║ 7} \\
\texttt{╚═══════╝} \\
\texttt{```} \\
            }
        }
    }
    & & \\ \\

    \theutterance \stepcounter{utterance}  
    & & \multicolumn{4}{p{0.6\linewidth}}{
        \cellcolor[rgb]{0.9,0.9,0.9}{
            \makecell[{{p{\linewidth}}}]{
                \texttt{\tiny{[P1$\langle$GM]}}
                \texttt{Il tuo messaggio è stato inoltrato all'altro giocatore.} \\
\\ 
\\ 
\texttt{Stai attualmente giocando il round 3 di massimo 12.} \\
\texttt{Avete accumulato collettivamente 0 di 8 penalità.} \\
\texttt{L'altro giocatore ha spostato un oggetto sulla sua griglia.} \\
\\ 
\\ 
\texttt{Qual è il tuo prossimo comando?} \\
            }
        }
    }
    & \\ \\

    \theutterance \stepcounter{utterance}  
    & \multicolumn{4}{p{0.6\linewidth}}{
        \cellcolor[rgb]{0.9,0.9,0.9}{
            \makecell[{{p{\linewidth}}}]{
                \texttt{\tiny{[P1$\rangle$GM]}}
                \texttt{COMUNICA: Tutto ok. Resto fermo; procedi e avvisami con "finito?" quando hai completato.} \\
            }
        }
    }
    & & \\ \\

    \theutterance \stepcounter{utterance}  
    & & & \multicolumn{2}{p{0.3\linewidth}}{
        \cellcolor[rgb]{0.9,0.9,0.9}{
            \makecell[{{p{\linewidth}}}]{
                \texttt{\tiny{[GM$|$GM]}}
                \texttt{COMUNICA: Tutto ok. Resto fermo; procedi e avvisami con "finito?" quando hai completato.} \\
            }
        }
    }
    & & \\ \\

    \theutterance \stepcounter{utterance}  
    & & \multicolumn{4}{p{0.6\linewidth}}{
        \cellcolor[rgb]{0.9,0.9,0.9}{
            \makecell[{{p{\linewidth}}}]{
                \texttt{\tiny{[GM$\rangle$P2]}}
                \texttt{Spostato 'L' in (5,1) con successo. La tua griglia aggiornata appare così:} \\
\texttt{```} \\
\texttt{ 1234567} \\
\texttt{╔═╤═══╤═╗} \\
\texttt{║◌│◌CL│◌║ 1} \\
\texttt{╟─┼──┬┘◌║ 2} \\
\texttt{║◌│◌P│◌◌║ 3} \\
\texttt{╟─┴──┘◌◌║ 4} \\
\texttt{║◌◌◌◌◌◌◌║ 5} \\
\texttt{║◌◌◌◌◌◌◌║ 6} \\
\texttt{║◌◌◌◌◌◌◌║ 7} \\
\texttt{╚═══════╝} \\
\texttt{```} \\
\texttt{Stai attualmente giocando il round 4 di massimo 12.} \\
\texttt{Avete accumulato collettivamente 0 di 8 penalità.} \\
\texttt{L'altro giocatore ha inviato questo messaggio:} \\
\\ 
\\ 
\texttt{"Tutto ok. Resto fermo; procedi e avvisami con "finito?" quando hai completato."} \\
\\ 
\\ 
\texttt{Qual è il tuo prossimo comando?} \\
            }
        }
    }
    & \\ \\

    \theutterance \stepcounter{utterance}  
    & & & \multicolumn{4}{p{0.6\linewidth}}{
        \cellcolor[rgb]{0.9,0.9,0.9}{
            \makecell[{{p{\linewidth}}}]{
                \texttt{\tiny{[GM$\langle$P2]}}
                \texttt{SPOSTA: C, (6, 5)} \\
            }
        }
    }
     \\ \\

    \theutterance \stepcounter{utterance}  
    & & & \multicolumn{2}{p{0.3\linewidth}}{
        \cellcolor[rgb]{0.9,0.9,0.9}{
            \makecell[{{p{\linewidth}}}]{
                \texttt{\tiny{[GM$|$GM]}}
                \texttt{SPOSTA: C, (6, 5)} \\
            }
        }
    }
    & & \\ \\

    \theutterance \stepcounter{utterance}  
    & & & \multicolumn{2}{p{0.3\linewidth}}{
        \cellcolor[rgb]{0.9,0.9,0.9}{
            \makecell[{{p{\linewidth}}}]{
                \texttt{\tiny{[GM$|$GM]}}
                \texttt{Spostato 'C' in (6,5) con successo. La tua griglia aggiornata appare così:} \\
\texttt{```} \\
\texttt{ 1234567} \\
\texttt{╔═╤═══╤═╗} \\
\texttt{║◌│◌◌L│◌║ 1} \\
\texttt{╟─┼──┬┘◌║ 2} \\
\texttt{║◌│◌P│◌◌║ 3} \\
\texttt{╟─┴──┘◌◌║ 4} \\
\texttt{║◌◌◌◌◌C◌║ 5} \\
\texttt{║◌◌◌◌◌◌◌║ 6} \\
\texttt{║◌◌◌◌◌◌◌║ 7} \\
\texttt{╚═══════╝} \\
\texttt{```} \\
            }
        }
    }
    & & \\ \\

    \theutterance \stepcounter{utterance}  
    & & \multicolumn{4}{p{0.6\linewidth}}{
        \cellcolor[rgb]{0.9,0.9,0.9}{
            \makecell[{{p{\linewidth}}}]{
                \texttt{\tiny{[P1$\langle$GM]}}
                \texttt{Il tuo messaggio è stato inoltrato all'altro giocatore.} \\
\\ 
\\ 
\texttt{Stai attualmente giocando il round 4 di massimo 12.} \\
\texttt{Avete accumulato collettivamente 0 di 8 penalità.} \\
\texttt{L'altro giocatore ha spostato un oggetto sulla sua griglia.} \\
\\ 
\\ 
\texttt{Qual è il tuo prossimo comando?} \\
            }
        }
    }
    & \\ \\

    \theutterance \stepcounter{utterance}  
    & \multicolumn{4}{p{0.6\linewidth}}{
        \cellcolor[rgb]{0.9,0.9,0.9}{
            \makecell[{{p{\linewidth}}}]{
                \texttt{\tiny{[P1$\rangle$GM]}}
                \texttt{COMUNICA: finito?} \\
            }
        }
    }
    & & \\ \\

    \theutterance \stepcounter{utterance}  
    & & & \multicolumn{2}{p{0.3\linewidth}}{
        \cellcolor[rgb]{0.9,0.9,0.9}{
            \makecell[{{p{\linewidth}}}]{
                \texttt{\tiny{[GM$|$GM]}}
                \texttt{COMUNICA: finito?} \\
            }
        }
    }
    & & \\ \\

    \theutterance \stepcounter{utterance}  
    & & \multicolumn{4}{p{0.6\linewidth}}{
        \cellcolor[rgb]{0.9,0.9,0.9}{
            \makecell[{{p{\linewidth}}}]{
                \texttt{\tiny{[GM$\rangle$P2]}}
                \texttt{Spostato 'C' in (6,5) con successo. La tua griglia aggiornata appare così:} \\
\texttt{```} \\
\texttt{ 1234567} \\
\texttt{╔═╤═══╤═╗} \\
\texttt{║◌│◌◌L│◌║ 1} \\
\texttt{╟─┼──┬┘◌║ 2} \\
\texttt{║◌│◌P│◌◌║ 3} \\
\texttt{╟─┴──┘◌◌║ 4} \\
\texttt{║◌◌◌◌◌C◌║ 5} \\
\texttt{║◌◌◌◌◌◌◌║ 6} \\
\texttt{║◌◌◌◌◌◌◌║ 7} \\
\texttt{╚═══════╝} \\
\texttt{```} \\
\texttt{Stai attualmente giocando il round 5 di massimo 12.} \\
\texttt{Avete accumulato collettivamente 0 di 8 penalità.} \\
\texttt{L'altro giocatore ha inviato questo messaggio:} \\
\\ 
\\ 
\texttt{"finito?"} \\
\\ 
\\ 
\texttt{Qual è il tuo prossimo comando?} \\
            }
        }
    }
    & \\ \\

    \theutterance \stepcounter{utterance}  
    & & & \multicolumn{4}{p{0.6\linewidth}}{
        \cellcolor[rgb]{0.9,0.9,0.9}{
            \makecell[{{p{\linewidth}}}]{
                \texttt{\tiny{[GM$\langle$P2]}}
                \texttt{COMUNICA: finito!} \\
            }
        }
    }
     \\ \\

    \theutterance \stepcounter{utterance}  
    & & & \multicolumn{2}{p{0.3\linewidth}}{
        \cellcolor[rgb]{0.9,0.9,0.9}{
            \makecell[{{p{\linewidth}}}]{
                \texttt{\tiny{[GM$|$GM]}}
                \texttt{COMUNICA: finito!} \\
            }
        }
    }
    & & \\ \\

    \theutterance \stepcounter{utterance}  
    & & & \multicolumn{2}{p{0.3\linewidth}}{
        \cellcolor[rgb]{0.9,0.9,0.9}{
            \makecell[{{p{\linewidth}}}]{
                \texttt{\tiny{[GM$|$GM]}}
                \texttt{true} \\
            }
        }
    }
    & & \\ \\

    \theutterance \stepcounter{utterance}  
    & & & \multicolumn{2}{p{0.3\linewidth}}{
        \cellcolor[rgb]{0.9,0.9,0.9}{
            \makecell[{{p{\linewidth}}}]{
                \texttt{\tiny{[GM$|$GM]}}
                \texttt{* success: True} \\
\texttt{* lose: False} \\
\texttt{* aborted: False} \\
\texttt{{-}{-}{-}{-}{-}{-}{-}} \\
\texttt{* Shifts: 2.00} \\
\texttt{* Max Shifts: 4.00} \\
\texttt{* Min Shifts: 2.00} \\
\texttt{* End Distance Sum: 0.00} \\
\texttt{* Init Distance Sum: 14.80} \\
\texttt{* Expected Distance Sum: 12.57} \\
\texttt{* Penalties: 0.00} \\
\texttt{* Max Penalties: 8.00} \\
\texttt{* Rounds: 5.00} \\
\texttt{* Max Rounds: 12.00} \\
\texttt{* Object Count: 3.00} \\
            }
        }
    }
    & & \\ \\

    \theutterance \stepcounter{utterance}  
    & & & \multicolumn{2}{p{0.3\linewidth}}{
        \cellcolor[rgb]{0.9,0.9,0.9}{
            \makecell[{{p{\linewidth}}}]{
                \texttt{\tiny{[GM$|$GM]}}
                \texttt{Initial Grids:} \\
\texttt{```} \\
\texttt{Player 1:     Player 2:} \\
\texttt{ 1234567       1234567} \\
\texttt{╔═╤═══╤═╗     ╔═╤═══╤═╗} \\
\texttt{║◌│◌◌L│◌║ 1   ║◌│◌C◌│◌║ 1} \\
\texttt{╟─┼──┬┘◌║ 2   ╟─┼──┬┘◌║ 2} \\
\texttt{║◌│◌P│◌◌║ 3   ║◌│◌◌│◌◌║ 3} \\
\texttt{╟─┴──┘◌◌║ 4   ╟─┴──┘◌◌║ 4} \\
\texttt{║◌◌◌◌◌C◌║ 5   ║◌◌◌◌◌◌◌║ 5} \\
\texttt{║◌◌◌◌◌◌◌║ 6   ║◌◌◌◌◌◌◌║ 6} \\
\texttt{║◌◌◌◌◌◌◌║ 7   ║◌◌◌P◌◌L║ 7} \\
\texttt{╚═══════╝     ╚═══════╝} \\
\texttt{```} \\
\\ 
\texttt{Final Grids:} \\
\texttt{```} \\
\texttt{Player 1:     Player 2:} \\
\texttt{ 1234567       1234567} \\
\texttt{╔═╤═══╤═╗     ╔═╤═══╤═╗} \\
\texttt{║◌│◌◌L│◌║ 1   ║◌│◌◌L│◌║ 1} \\
\texttt{╟─┼──┬┘◌║ 2   ╟─┼──┬┘◌║ 2} \\
\texttt{║◌│◌P│◌◌║ 3   ║◌│◌P│◌◌║ 3} \\
\texttt{╟─┴──┘◌◌║ 4   ╟─┴──┘◌◌║ 4} \\
\texttt{║◌◌◌◌◌C◌║ 5   ║◌◌◌◌◌C◌║ 5} \\
\texttt{║◌◌◌◌◌◌◌║ 6   ║◌◌◌◌◌◌◌║ 6} \\
\texttt{║◌◌◌◌◌◌◌║ 7   ║◌◌◌◌◌◌◌║ 7} \\
\texttt{╚═══════╝     ╚═══════╝} \\
\texttt{```} \\
            }
        }
    }
    & & \\ \\

\end{supertabular}
}

\end{document}
